\documentclass[12pt]{article}

\usepackage{sbc-template}
\usepackage[brazil]{babel}
\usepackage[utf8]{inputenc}
\usepackage[T1]{fontenc}
\usepackage{graphicx}
\usepackage{url}
\usepackage{hyperref}

\sloppy

\title{Migração de um Sistema de Apoio à Decisão em Saúde do\\
Paradigma PWA para Aplicativo Móvel: Uma Análise de\\
Engenharia de Software}

\author{Ana Carolina Dias Faria\inst{1}, Diogo S. Martins\inst{1}}

\address{Centro de Matemática, Computação e Cognição\\
  Universidade Federal do ABC (UFABC) -- Santo André -- SP -- Brasil
  \email{ana.faria@aluno.ufabc.edu.br, diogo.martins@ufabc.edu.br}
}

\begin{document}

\maketitle

%-----------------------------------------------------------------------
\begin{abstract}
This paper investigates the architectural paradigm shift in digital health
by analyzing the migration of a Clinical Decision Support System (CDSS)
from a Progressive Web App (PWA) to a dedicated mobile platform.
Built on an API-first approach using FastAPI, the system performs text
parsing and OCR on blood test reports to provide patient-centered
explanations. While PWAs offer immediate access, web-view-based
frameworks impose significant CPU and energy overhead during basic
interactions. This case study evaluates the migration to React Native,
characterizing software engineering benefits such as perceived
performance, architectural reuse, and usability heuristics.
\end{abstract}

\begin{resumo}
Este artigo investiga a mudança de paradigma arquitetural em saúde
digital por meio da análise da migração de um Sistema de Apoio à
Decisão Clínica (SADC) de uma aplicação web progressiva (PWA) para
uma plataforma móvel dedicada. Estruturado em uma abordagem API-first
com FastAPI, o sistema realiza \textit{parsing} de texto e OCR em laudos
de hemogramas para fornecer \textit{briefings} educativos aos pacientes.
Embora PWAs ofereçam acesso imediato, \textit{frameworks} baseados em
\textit{web views} impõem \textit{overhead} severo de CPU e energia em
interações básicas. Este estudo de caso avalia a migração para React
Native, caracterizando empiricamente os benefícios de engenharia de
software em termos de desempenho percebido, reuso arquitetural e
heurísticas de usabilidade.
\end{resumo}

%-----------------------------------------------------------------------
\section{Introdução}

O hemograma é um dos exames laboratoriais mais requisitados no Brasil,
mas sua interpretação por usuários comuns é difícil, sendo agravada pelo
fato de os valores de referência impressos nos laudos derivarem, em geral,
de populações estrangeiras \cite{rosenfeld2019}. O projeto-base deste
trabalho, denominado InterpreteLabBR, foi concebido originalmente na
graduação como uma Progressive Web App (PWA). O sistema atua na extração
de valores de hemogramas no padrão do SUS por meio de \textit{parsing} de
texto e reconhecimento óptico de caracteres (OCR), classificando cada
analito segundo as faixas estabelecidas pela Pesquisa Nacional de Saúde
\cite{rosenfeld2019} e gerando um \textit{briefing} educativo ao paciente.

Embora as PWAs ofereçam a vantagem de um acesso imediato, a literatura
aponta limitações críticas de desempenho frente a soluções móveis
dedicadas. \cite{oliveira2023} demonstram que \textit{frameworks} baseados
em \textit{web views} impõem um \textit{overhead} severo de CPU e energia
em interações básicas, como a rolagem de tela, chegando a demandar até 8
vezes mais energia do que soluções estruturadas sobre componentes nativos.
Diante desse cenário, a motivação desta pesquisa é investigar e
documentar, sob a ótica da Engenharia de Software, os impactos técnicos e
operacionais de migrar esse artefato de um PWA encapsulado para uma
arquitetura móvel dedicada \cite{oliveira2023}.

A relevância social e científica desta transição reflete-se diretamente
nos seus potenciais beneficiários. Por um lado, os pacientes usuários do
SUS ganham uma interface mais responsiva e um canal de distribuição mais
inclusivo por meio de um aplicativo instalável, alinhando-se ao eixo de
cidadania e inclusão apontado pelo Brazilian Digital Health Index (BDHI)
\cite{cruz2022}. Por outro lado, a comunidade de Engenharia de Software
aplicada à saúde beneficia-se com um estudo de caso empírico e rico em
métricas computacionais, detalhando os benefícios, os custos e os desafios
práticos de migrar um sistema de informação em saúde entre diferentes
paradigmas de distribuição \cite{oliveira2023}.

A partir desse contexto, estabelece-se o problema de pesquisa deste
trabalho, que busca compreender quais são, e como podem ser medidos de
forma reprodutível, os benefícios de Engenharia de Software especificamente
em termos de desempenho percebido, reuso arquitetural e usabilidade
decorrentes da migração de um sistema de apoio à decisão clínica do
paradigma PWA para um aplicativo móvel dedicado, garantindo a preservação
da integridade funcional do sistema \cite{rampinelli2026}.

Ao analisar como esse problema é tratado atualmente na literatura,
identificam-se frentes de soluções existentes que evidenciam a necessidade
deste estudo. As ferramentas comerciais de leitura de exames operam
majoritariamente em nuvem, dependem de parâmetros populacionais
estrangeiros e desconsideram tanto o leiaute de impressão do SUS quanto a
referência epidemiológica nacional. Similarmente, os leitores de OCR
comerciais genéricos limitam-se a extrair o texto bruto do documento, sem
realizar qualquer tipo de interpretação clínica ou contextualização para o
entendimento do paciente leigo. Por fim, o próprio protótipo inicial do
InterpreteLabBR em PWA, embora funcional, carece de uma validação formal e
sofre com as restrições de encapsulamento e de acesso a \textit{hardware}
impostas pelos navegadores \cite{cardieri2018}. As soluções comerciais
vigentes funcionam, em geral, como ``caixas-pretas'' inadequadas à
realidade técnica do sistema público de saúde. Portanto, persiste na
literatura uma lacuna de estudos de caso reais que quantifiquem e mapeiem
de forma clara os efeitos práticos de uma migração de paradigma arquitetural
sobre sistemas de informação em saúde, lacuna esta que este trabalho se
propõe a preencher \cite{rampinelli2026}.

%-----------------------------------------------------------------------
\section{Revisão de Literatura}

\subsection{Sistemas de Apoio à Decisão Clínica na Saúde Digital Brasileira}

Os Sistemas de Apoio à Decisão Clínica (\textit{Clinical Decision Support
Systems} — CDSS) são ferramentas computacionais projetadas para auxiliar
profissionais de saúde e pacientes no processo de tomada de decisão,
fornecendo informações baseadas em evidências em contextos clínicos
específicos. De acordo com a literatura, essas ferramentas integram dados
do paciente com bases de conhecimento estruturadas para gerar
recomendações, alertas ou interpretações que auxiliam na triagem e no
diagnóstico \cite{bauermann2024}.

No contexto brasileiro, a saúde digital tem sido avaliada por iniciativas
como o Brazilian Digital Health Index (BDHI), instrumento desenvolvido para
mensurar o nível de maturidade da saúde digital no país. \cite{cruz2022}
demonstram que o Brasil apresenta avanços significativos em alguns eixos,
como conectividade e infraestrutura, mas ainda enfrenta lacunas expressivas
no eixo de cidadania e inclusão, que abrange o acesso da população a
ferramentas digitais de saúde. Esse diagnóstico evidencia a necessidade de
soluções que priorizem a acessibilidade e a inclusão digital como critério
de projeto, e não apenas como efeito secundário.

No campo dos CDSS aplicados à saúde pública brasileira, destaca-se o
trabalho de \cite{bauermann2024}, que desenvolveram um sistema de apoio à
decisão clínica voltado a cuidadores de pacientes pediátricos hemofílicos.
Esse trabalho exemplifica como sistemas computacionais podem ser projetados
para mediar o conhecimento clínico especializado e o leigo, reduzindo
barreiras de interpretação para familiares e cuidadores sem formação na
área da saúde, em uma proposta análoga à do InterpreteLabBR, que busca
traduzir resultados laboratoriais de hemogramas para o paciente comum do
SUS. A relevância de soluções móveis nesse contexto é reforçada por
\cite{torrente2024}, que analisam a relação entre tecnologia e atendimento
pré-hospitalar móvel, evidenciando que a adoção de aplicativos instaláveis
no contexto de saúde pública contribui para a agilidade e a qualidade do
atendimento. A disponibilização de sistemas de saúde em plataformas móveis
instaláveis configura, portanto, não apenas uma decisão técnica, mas uma
estratégia de alcance e equidade no acesso à informação clínica.

\subsection{Hemograma e Valores de Referência Populacionais Brasileiros}

O hemograma é um dos exames laboratoriais mais solicitados na rotina
clínica do Sistema Único de Saúde (SUS), permitindo a avaliação de
componentes do sangue como eritrócitos, leucócitos e plaquetas. Apesar
de sua ampla utilização, a interpretação dos resultados por pacientes
leigos é frequentemente dificultada pela ausência de contextualização
adequada nos laudos e, sobretudo, pelo fato de que os valores de referência
impressos na maioria dos laudos brasileiros derivam de estudos populacionais
estrangeiros, que não refletem as especificidades demográficas, nutricionais
e epidemiológicas da população brasileira.

Essa lacuna foi endereçada de forma expressiva pela Pesquisa Nacional de
Saúde (PNS), cujos resultados para exames laboratoriais de hemograma foram
publicados por \cite{rosenfeld2019}. O estudo, de abrangência nacional,
estabeleceu valores de referência para a população adulta brasileira
estratificados por sexo e faixa etária, constituindo o primeiro
levantamento paramétrico de grande escala conduzido com amostra
representativa do Brasil. A adoção desses parâmetros como base
classificatória confere ao InterpreteLabBR validade epidemiológica superior
às ferramentas que utilizam tabelas de referência importadas, tornando-o
mais adequado à realidade do sistema público de saúde nacional. A
importância de se utilizar valores de referência populacionalmente ajustados
reside no fato de que a classificação equivocada de um analito como anormal,
ou, inversamente, a não detecção de um valor limítrofe, pode gerar ansiedade
desnecessária no paciente ou, em casos mais graves, retardar a busca por
atendimento médico. A fundamentação do sistema na PNS representa, portanto,
um diferencial de segurança e de relevância clínica para a população usuária
do SUS, público-alvo primário do projeto.

\subsection{\textit{Progressive Web Apps}: Características e Limitações}

As \textit{Progressive Web Apps} (PWAs) constituem uma abordagem de
desenvolvimento de software que busca combinar as vantagens das aplicações
web — acessibilidade via navegador, sem necessidade de instalação — com
características típicas de aplicativos nativos, como comportamento
responsivo, suporte a notificações e capacidade de funcionamento
\textit{offline} por meio de \textit{Service Workers}.

No contexto do InterpreteLabBR, a escolha inicial pela arquitetura PWA
justificou-se pela praticidade de distribuição e pela dispensa de contas em
plataformas comerciais. No entanto, essa abordagem impõe restrições
estruturais que limitam o desempenho e a experiência do usuário. As PWAs
executam sobre o motor de renderização do navegador (\textit{web view}), o
que introduz uma camada de abstração entre o código da aplicação e os
componentes nativos do sistema operacional. Essa camada adicional implica
\textit{overhead} de processamento, especialmente em operações que
demandam alta responsividade de interface, como rolagem de listas extensas,
animações ou interações rápidas com o teclado. Adicionalmente, PWAs
enfrentam limitações de integração com recursos de \textit{hardware} do
dispositivo de forma mais restrita do que aplicativos nativos, por estarem
sujeitas às permissões e às políticas de segurança impostas pelo navegador.
No contexto do InterpreteLabBR, onde o carregamento de documentos PDF via
seletor de arquivos nativo é uma operação central, essas restrições
impactam diretamente a experiência do usuário e motivam a evolução para uma
arquitetura móvel dedicada.

\subsection{\textit{Frameworks} de Desenvolvimento Móvel e \textit{Overhead} de Recursos}

O desenvolvimento de aplicativos móveis pode ser realizado por meio de três
abordagens principais: (i) desenvolvimento nativo, em que a aplicação é
escrita nas linguagens oficiais da plataforma, com acesso irrestrito às
APIs do sistema operacional; (ii) desenvolvimento híbrido baseado em
\textit{web view}, em que o código web é encapsulado em um contêiner nativo
e renderizado por um motor de navegador embutido; e (iii) \textit{frameworks}
de renderização nativa, como o React Native, que traduzem componentes
escritos em JavaScript/TypeScript em componentes nativos da plataforma, sem
o uso de \textit{web views}.

A distinção entre essas abordagens tem implicações diretas no desempenho e
no consumo de recursos do dispositivo. \cite{oliveira2023} conduziram um
estudo empírico sobre o \textit{overhead} de recursos em
\textit{frameworks} de desenvolvimento móvel, comparando soluções baseadas
em \textit{web views} com aquelas baseadas em componentes nativos. Os
resultados demonstram que \textit{frameworks} baseados em \textit{web views}
podem demandar até 8 vezes mais energia do que soluções estruturadas sobre
componentes nativos em interações básicas, como a rolagem de listas. Esse
\textit{overhead} impacta diretamente a experiência do usuário em
dispositivos com recursos limitados de bateria e processamento, cenário
comum entre usuários do SUS, que frequentemente utilizam aparelhos Android
de entrada. O React Native, \textit{framework} adotado no InterpreteLabBR
para a camada móvel, posiciona-se como uma solução intermediária que combina
a produtividade do desenvolvimento em JavaScript/TypeScript com a eficiência
de renderização de componentes nativos. Ao contrário de soluções que
encapsulam \textit{web views}, o React Native mapeia os componentes
declarativos da aplicação diretamente para os componentes visuais do sistema
operacional, favorecendo uma execução mais fluida e um menor consumo de
recursos. O Expo fornece uma camada de abstração que simplifica o processo
de \textit{build} e distribuição, incluindo a geração de APKs por meio do
serviço EAS Build sem a necessidade de ambiente nativo configurado
localmente.

\subsection{Arquitetura \textit{API-first} e Reuso Arquitetural em Migrações de Software}

A arquitetura \textit{API-first} é uma abordagem de \textit{design} de
software na qual a interface de programação de aplicações (API) é concebida
e documentada antes da implementação dos clientes que a consomem. Nesse
modelo, toda a lógica de negócio e processamento de dados é centralizada no
servidor (\textit{backend}), exposta por meio de \textit{endpoints} REST ou
GraphQL, e consumida por diferentes clientes de forma independente. Essa
separação de responsabilidades constitui um princípio fundamental de
Engenharia de Software, alinhado ao conceito de baixo acoplamento e alta
coesão.

No contexto de migrações de sistema, a arquitetura \textit{API-first}
apresenta vantagem expressiva: ao concentrar a lógica clínica no
\textit{backend}, ela permite que a substituição ou a evolução da camada de
apresentação — de uma PWA para um aplicativo móvel — ocorra sem que o
núcleo funcional do sistema precise ser reescrito. No caso do
InterpreteLabBR, todo o processamento crítico, incluindo a extração de
valores via OCR e a classificação dos analitos segundo as faixas da PNS,
reside no \textit{backend} implementado em FastAPI (Python), exposto como
serviço REST. Isso permite que tanto o cliente web (PWA) quanto o cliente
móvel (React Native) consumam os mesmos \textit{endpoints}, garantindo
consistência nos resultados e eliminando a necessidade de duplicação de
lógica clínica. Esse padrão arquitetural é também relevante do ponto de
vista da manutenibilidade e da evolução do sistema. \cite{rampinelli2026}
propõem um \textit{framework} de avaliação de sistemas de informação sob a
perspectiva do usuário, destacando que sistemas bem estruturados
arquiteturalmente tendem a apresentar maior facilidade de adaptação a novos
contextos de uso e plataformas. A documentação e a caracterização formal do
reuso arquitetural constituem, portanto, uma contribuição relevante deste
trabalho para a literatura de Engenharia de Software aplicada à saúde.

\subsection{Reconhecimento Óptico de Caracteres em Documentos de Saúde}

O Reconhecimento Óptico de Caracteres (\textit{Optical Character Recognition}
— OCR) é uma tecnologia que converte imagens de texto em texto digital
estruturado, possibilitando o processamento computacional de documentos
digitalizados ou gerados por impressão. No contexto de sistemas de saúde,
o OCR viabiliza a extração automatizada de dados a partir de laudos,
prescrições e prontuários, reduzindo a necessidade de entrada manual de
informações e ampliando a escalabilidade de soluções de triagem e análise
clínica.

A qualidade da extração por OCR em documentos de saúde está sujeita a
desafios específicos, incluindo variações no \textit{layout} de impressão,
presença de ruídos gráficos, diferenças tipográficas entre unidades de saúde
e inconsistências na qualidade de digitalização. \cite{most2025} abordam
abordagens baseadas em visão computacional para recuperação robusta de
documentos, evidenciando que estratégias que combinam OCR com técnicas de
pré-processamento de imagem tendem a apresentar maior resiliência a variações
de qualidade. No contexto do InterpreteLabBR, o leiaute padronizado do laudo
do SUS representa tanto uma vantagem — ao permitir o desenvolvimento de
heurísticas de \textit{parsing} específicas para esse formato — quanto um
desafio, dado que pequenas variações entre laboratórios da rede pública podem
introduzir erros de extração.

\cite{lima2024} comparou modelos de OCR de código aberto --- Tesseract e
EasyOCR --- com a API Google Cloud Vision para extração de informações em
prescrições médicas, avaliando precisão, velocidade e robustez em diferentes
cenários de caligrafia e condições de digitalização. Os resultados
demonstraram que modelos de OCR tradicionais apresentam vantagens em
velocidade e precisão para textos impressos, enquanto modelos especializados
em reconhecimento de texto manuscrito superam as demais opções em caligrafias
variadas \cite{lima2024}. Esse achado é diretamente relevante para o
InterpreteLabBR, que adota o Tesseract como estratégia de \textit{fallback}
na extração de laudos hematológicos do SUS: laudos gerados por sistemas
informatizados são processados com maior confiabilidade pelos mesmos modelos
de OCR open-source avaliados por Lima, enquanto documentos digitalizados por
câmera ou scanner de baixa qualidade exigem estratégias adicionais de
pré-processamento de imagem. A questão do
pós-processamento de texto extraído por OCR em língua portuguesa é
endereçada por \cite{osorio2025}, que propõem recursos específicos para
otimização de texto em português após extração por OCR. Essa contribuição é
particularmente relevante para o contexto do InterpreteLabBR, cujos laudos
são emitidos em português e podem conter termos técnicos da nomenclatura
hematológica que ferramentas genéricas de OCR tendem a interpretar de forma
equivocada.

\subsection{Avaliação de Usabilidade por Inspeção Heurística}

A usabilidade é um atributo de qualidade de software que mede a facilidade
com que usuários conseguem aprender e utilizar uma interface para alcançar
seus objetivos, bem como o grau de satisfação proporcionado durante o uso.
No contexto de sistemas de saúde voltados ao paciente leigo, a usabilidade
assume dimensão crítica: interfaces confusas ou que exigem conhecimento
técnico prévio podem resultar em interpretações equivocadas dos dados
clínicos, comprometendo o propósito informativo do sistema.

Entre os métodos de avaliação de usabilidade, a Avaliação Heurística,
proposta por Jakob Nielsen, é amplamente utilizada na prática de Engenharia
de Software por sua eficiência e pelo baixo custo de aplicação. O método
consiste na inspeção sistemática da interface por especialistas em
usabilidade ou interação humano-computador (IHC), que avaliam a adequação
da interface a um conjunto de princípios estabelecidos — as heurísticas de
Nielsen — registrando e classificando os problemas encontrados segundo uma
escala de severidade de 0 a 4 \cite{cardieri2018}. A ausência de usuários
reais no processo de avaliação, substituídos por especialistas que atuam
como juízes inspetores, é uma característica que distingue a avaliação
heurística dos testes de usabilidade com usuários, dispensando aprovação de
comitê de ética e viabilizando a execução dentro de prazos acadêmicos mais
curtos \cite{cardieri2018}.

\cite{cardieri2018} realizaram um estudo comparativo da experiência do
usuário em aplicações \textit{web} móveis, nativas e PWAs, sob as
perspectivas de usuários reais e especialistas em IHC. Os resultados
evidenciam diferenças perceptíveis de qualidade de interação entre os
paradigmas, com aplicações nativas e PWAs bem implementadas tendendo a
superar aplicações \textit{web} móveis em critérios como responsividade,
fluidez e consistência de navegação. Esse achado fundamenta a hipótese
central deste trabalho, de que a migração do InterpreteLabBR para um
aplicativo móvel resultará em ganhos mensuráveis de usabilidade,
verificáveis pela inspeção heurística comparativa. \cite{rampinelli2026}
propõem um \textit{framework} de avaliação de sistemas de informação sob a
perspectiva do usuário, reforçando a necessidade de metodologias estruturadas
para a avaliação da qualidade percebida em sistemas de saúde digital.
\cite{vaz2026} desenvolveu e avaliou a usabilidade do APScreening, um
aplicativo móvel de apoio à decisão clínica para rastreamentos na Atenção
Primária à Saúde no contexto do SUS, obtendo pontuação SUS de 91,25 ---
classificada como excelente ---, evidenciando que CDSSs móveis voltados ao
ecossistema de saúde pública brasileiro podem atingir alta aceitabilidade em
condições reais de uso.

\subsection{\textit{Design Science Research} como Abordagem Metodológica}

A \textit{Design Science Research} (DSR) é uma abordagem metodológica
originada na área de Sistemas de Informação, utilizada para orientar a
concepção, o desenvolvimento e a avaliação de artefatos que resolvem
problemas práticos identificados. Diferentemente das metodologias de
pesquisa baseadas em observação de fenômenos existentes, a DSR parte da
construção de um artefato — seja um sistema, um modelo, um método ou um
constructo — e avalia sua efetividade no contexto do problema que motivou
sua criação.

A DSR é especialmente adequada a trabalhos de pós-graduação em Tecnologia e
Sistemas de Informação que envolvem o desenvolvimento de soluções
computacionais, pois confere rigor científico ao processo de construção e
validação do artefato, exigindo que o problema seja claramente identificado,
que a solução seja fundamentada na literatura e que sua avaliação seja
conduzida de forma sistemática e reprodutível. No caso do InterpreteLabBR,
o artefato é o aplicativo móvel resultante da migração do PWA original, e
sua avaliação ocorre nas três frentes metodológicas previstas: desempenho
comparativo (Frente A), caracterização da engenharia da migração (Frente B)
e inspeção de usabilidade (Frente C).

%-----------------------------------------------------------------------
\section{Trabalhos Relacionados}

A presente seção situa o InterpreteLabBR em relação a trabalhos recentes
que investigam desempenho, migração e desafios de desenvolvimento em
\textit{frameworks} de aplicativos móveis \textit{cross-platform},
identificando as contribuições específicas que o diferenciam da literatura
existente.

\cite{sokolov2025} conduziu um estudo experimental quantitativo comparando
o consumo de memória e a eficiência de tempo de execução entre uma aplicação
React Native e seu equivalente nativo em Kotlin, utilizando criptografia e
decriptografia AES como carga computacional de referência e o Android Debug
Bridge (ADB) como instrumento de coleta de métricas. Os resultados
demonstraram que a implementação nativa superou o React Native em eficiência
de CPU e crescimento de memória ao longo do tempo, embora em determinados
cenários de decriptografia o React Native tenha apresentado desempenho
superior. O trabalho contribui com evidências empíricas sobre o
\textit{overhead} de execução do React Native em tarefas computacionalmente
intensivas, mas restringe-se a um \textit{benchmark} sintético — criptografia
de arquivos — sem qualquer vínculo com aplicações de domínio específico ou
com o processo de migração de sistemas existentes. O InterpreteLabBR
complementa essa perspectiva ao avaliar o \textit{overhead} de desempenho
em um cenário real de uso, com operações predominantemente
\textit{I/O-bound} e \textit{rede-bound} (submissão de PDF e retorno de
JSON da API), sobre dispositivos Android de entrada representativos do
perfil de usuários do SUS.

\cite{zhou2024} realizou um estudo qualitativo com 20 desenvolvedores
experientes em React Native e Flutter, por meio de entrevistas
semiestruturadas, para mapear os desafios e soluções em projetos de
desenvolvimento \textit{cross-platform}. Entre os desafios mais recorrentes
identificados, destacam-se a otimização de desempenho de renderização, a
consistência da interface entre plataformas, a integração com módulos
nativos e a curva de aprendizado associada às atualizações frequentes dos
\textit{frameworks}. O trabalho oferece uma perspectiva valiosa sobre a
experiência do desenvolvedor (DX), mas opera exclusivamente no plano
qualitativo e não documenta o processo de migração de um sistema
preexistente de um paradigma para outro. O InterpreteLabBR avança nessa
direção ao registrar empiricamente o processo de migração de um PWA em
produção para React Native, quantificando o reuso de componentes,
identificando os padrões de portabilidade adotados e caracterizando o
esforço de adaptação por categoria de artefato, informação ausente na
literatura revisada por Zhou.

\cite{karin2025} conduziu uma revisão sistemática da literatura sobre o
impacto de \textit{frameworks cross-platform} no desempenho de aplicações
móveis, sintetizando evidências que posicionam soluções baseadas em
\textit{WebView} — incluindo PWAs — como as de maior \textit{overhead}
dentre as abordagens analisadas, com médias de 30 a 60 quadros por segundo
e consumo muito elevado de CPU, enquanto o React Native se mostrou mais
eficiente para aplicações orientadas a dados. A revisão conclui que, para
80\% das aplicações de negócio com perfil CRUD, a diferença de desempenho
entre paradigmas é imperceptível ao usuário final. Embora o trabalho
consolide evidências comparativas relevantes, trata-se de uma síntese da
literatura existente, sem coleta primária de dados e sem estudo de caso
aplicado a um sistema real. O InterpreteLabBR posiciona-se como um estudo
de caso primário que instancia, em um sistema de saúde pública real,
exatamente a transição de paradigma identificada por Karin como
teoricamente vantajosa — de PWA para React Native —, verificando
empiricamente se os ganhos de desempenho previstos pela literatura se
manifestam no contexto específico de laudos hematológicos no SUS.

Em síntese, os três trabalhos relacionados abordam, de formas
complementares, questões de desempenho e experiência de desenvolvimento em
\textit{frameworks cross-platform} \cite{sokolov2025,zhou2024,karin2025},
mas nenhum deles trata simultaneamente de: (i) um sistema real de apoio à
decisão clínica como objeto de avaliação; (ii) o processo de migração de
um PWA em produção para React Native como fenômeno de pesquisa; e (iii) a
avaliação tripartite que integra desempenho técnico, caracterização da
engenharia de migração e inspeção de usabilidade. A originalidade do
InterpreteLabBR reside precisamente nessa intersecção, ao conduzir um
estudo de caso empírico e reprodutível em um sistema de saúde pública
nacional, utilizando ferramentas gratuitas e metodologia verificável, que
possa servir de referência para migrações similares no contexto da saúde
digital brasileira.

%-----------------------------------------------------------------------
% Seções futuras (a preencher):
% \section{Arquitetura do Sistema e Processo de Migração}
% \section{Metodologia de Avaliação}
% \section{Resultados}
% \section{Conclusão}

%-----------------------------------------------------------------------
\bibliographystyle{sbc}
\bibliography{referencias}

\end{document}
